% swiftui-app-scenes-lifecycle.tex — the SwiftUI App protocol and @main, Scene types (WindowGroup,
% WindowGroup(for:), DocumentGroup, Window, Settings, MenuBarExtra), scenePhase at view vs app level,
% UIApplicationDelegateAdaptor + a scene delegate via configurationForConnecting, onOpenURL /
% onContinueUserActivity / handlesExternalEvents, openURL / openWindow / dismissWindow, @SceneStorage vs
% @AppStorage, backgroundTask(.appRefresh), commands, TN3187 (UIScene lifecycle required).
% Not repeated: the UIKit scene object model + restoration flow (ios-platform/state-restoration-multiwindow),
% app states / VC callbacks (ios-swift/app-and-vc-lifecycle), BGTaskScheduler internals
% (ios-platform/background-execution), deep-link routing into a path (swiftui-navigation).
% Sources (checked 2026-09-26): developer.apple.com/documentation/swiftui — ScenePhase (view =
% containing scene; App = aggregate, active if any scene is active), UIApplicationDelegateAdaptor (iOS 14;
% one per App, runtime error if declared twice; ObservableObject delegate goes into the environment;
% scene delegate via application(_:configurationForConnecting:options:) + delegateClass), Window (macOS 13,
% visionOS 26, not iOS), Settings (macOS 11), MenuBarExtra (macOS 13), DocumentGroup (iOS 14),
% OpenWindowAction (iOS 16; an existing window for the value is brought to front), DismissWindowAction
% (iOS 17), SceneStorage (lightweight, not for sensitive data, destroyed with the scene),
% onOpenURL(perform:) (universal links delivered as URLs; see handlesExternalEvents),
% onOpenURL(prefersInApp:) (26), Scene.backgroundTask(_:action:) (iOS 16; done when the action returns,
% then suspended; cancelled on timeout), Scene.commands(content:) (iPadOS: ⌘ HUD + menu bar);
% TN3187 "Migrating to the UIKit scene-based life cycle" (release after iOS 26: UIKit apps built with the
% latest SDK without UIScene life cycle will not launch).
% @labels: area=swiftui kind=api platform=apple topic=ui,platform-apis,architecture discussion=176
% @tags: app-protocol, main-attribute, windowgroup, documentgroup, scenephase, uiapplicationdelegateadaptor, onopenurl, openwindow, scenestorage, appstorage, backgroundtask, commands
\documentclass[8pt]{extarticle}
\usepackage{printup-sheet}
\usepackage{array}

\lstdefinelanguage{SwiftSheet}{
  morekeywords={func,let,var,struct,class,final,enum,return,if,else,guard,await,async,try,
    case,self,some,in,private,for,true,false,nil,
    @main,@State,@Environment,@SceneStorage,@AppStorage,@UIApplicationDelegateAdaptor,@MainActor},
  alsoletter={@}, sensitive=true, morecomment=[l]{//}, morestring=[b]",
  literate={->}{{\hbox{-}\hbox{>}}}2 {==}{{\hbox{=}\hbox{=}}}2 {??}{{\hbox{?}\hbox{?}}}2
           {!=}{{\hbox{!}\hbox{=}}}2 {...}{{\hbox{.}\hbox{.}\hbox{.}}}3}
\newcolumntype{L}[1]{>{\raggedright\arraybackslash}p{#1}}

\tikzset{
  lbl/.style={font=\scriptsize, text=black!80, inner sep=1pt, align=center},
  nd/.style={box, font=\scriptsize, inner sep=1.5pt, minimum height=5mm},
  sc/.style={nd, draw=sheetBrown, fill=sheetBrown!8},
  win/.style={nd, draw=sheetBlue, fill=sheetBlue!8, text width=19mm},
  ph/.style={font=\tiny\bfseries, rounded corners=1pt, inner sep=1.2pt},
}

\begin{document}

\sheettitle{SwiftUI App, scenes \& lifecycle}{swiftui · folio}

\oneliner{A SwiftUI app is a value: \texttt{@main struct MyApp: App} whose \texttt{body} returns
\textbf{scenes}. A scene is a \textbf{template} the system instantiates into one or more \textbf{windows}
(\texttt{UIWindowScene}s), each with its \textbf{own} view tree, \texttt{@State}, \texttt{@SceneStorage} and
\texttt{scenePhase}. UIKit's delegates come back through an \textbf{adaptor}; URLs, activities and background
wake-ups are delivered to scenes by \textbf{modifiers}.}

\begin{multicols}{2}
\raggedright

\section{How it works}
\begin{itemize}
  \item \textbf{\texttt{@main}} on an \texttt{App} generates \texttt{main()}. \texttt{App.init()} runs once,
        before any scene. \texttt{@State} on the App = process-wide model, injected into every window via
        \texttt{.environment(store)}; state in a root \emph{view} = one copy \textbf{per window}.
  \item \textbf{\texttt{scenePhase}} (\texttt{.active / .inactive / .background}): in a view = the
        \emph{containing} scene; in the \texttt{App} = \textbf{aggregate} — \texttt{.active} if any scene is,
        \texttt{.background} = app going to background. \texttt{.inactive} = visible, not interactive
        (switcher, Control Center) — pause, don't tear down.
  \item \textbf{UIKit hooks}: \texttt{@UIApplicationDelegateAdaptor} — once, in the \texttt{App} (twice =
        runtime error); SwiftUI creates it; an \texttt{ObservableObject} one lands in the environment. Push
        token, SDK setup. Scene delegate: \texttt{delegateClass} on the \texttt{UISceneConfiguration} returned
        from \texttt{configurationForConnecting}.
  \item \textbf{Incoming}: \texttt{.onOpenURL} — custom schemes \emph{and} Universal Links, as URLs;
        \texttt{.onContinueUserActivity(type)} — Handoff, Spotlight; \texttt{.handlesExternalEvents(matching:)}
        picks the window.
  \item \textbf{Outgoing} (environment): \texttt{openURL} (own \texttt{OpenURLAction}: \texttt{.handled /
        .discarded / .systemAction}); \texttt{openWindow(id:/value:)} (16) — a window already showing the value
        comes forward; \texttt{dismissWindow} (17); \texttt{supportsMultipleWindows}.
  \item \textbf{Persistence}: \texttt{@SceneStorage} — per window, system-restored, small, \textbf{no
        secrets}, dies with the scene. \texttt{@AppStorage} — \texttt{UserDefaults}, shared by all windows.
  \item \textbf{\texttt{.backgroundTask(.appRefresh("id"))}} (16): done when the closure \textbf{returns},
        then suspended; cancelled on timeout. Still needs the id in
        \texttt{BGTaskSchedulerPermittedIdentifiers} + a \emph{submitted} \texttt{BGAppRefreshTaskRequest}.
  \item \textbf{\texttt{.commands \{ CommandMenu \ldots \}}}: macOS menu bar; iPadOS: Cmd-hold HUD + menu bar.
\end{itemize}

\section{Scene types}
{\footnotesize\setlength{\tabcolsep}{3pt}
\begin{tabular}{@{}L{34mm}L{11mm}L{33mm}@{}}
\toprule
\texttt{WindowGroup} & all & template; many windows on iPad / Mac \\
\texttt{WindowGroup(for: T.self)} & 16 & a window per value \\
\texttt{DocumentGroup} & 14 & browser, open / save, window per doc \\
\texttt{Window} & macOS 13 & single unique window, not iOS \\
\texttt{Settings}, \texttt{MenuBarExtra} & macOS & preferences; menu-bar item \\
\bottomrule
\end{tabular}\par}

\section{Remember}
\textbf{``App = the process, Scene = a template, Window = an instance with its own state and phase.''}


\columnbreak

\section{App, scenes, windows, phases}
\begin{tikzpicture}[sheet]
  \node[nd, draw=sheetGreen!70!black, fill=sheetGreen!10, minimum width=74mm] (app) at (3.6,0)
       {\textbf{\texttt{NotesApp: App}} — one process · \texttt{@State store} · adaptor};
  \node[nd, draw=sheetGrey, fill=black!4] (ad) at (2.8,-0.75) {\texttt{AppDelegate}};
  \node[nd, draw=sheetGrey, fill=black!4] (sd) at (4.95,-0.75) {\texttt{SceneDelegate}};
  \draw[sheetGrey, dashed, ->] (ad) -- (sd);
  \node[lbl, text=black!60, font=\tiny] at (3.9,-1.12) {\texttt{delegateClass}};
  \draw[sheetGrey, dashed] (ad.north) -- (ad.north |- app.south);
  % scenes
  \node[sc] (s1) at (1.3,-1.6) {scene: \texttt{WindowGroup}};
  \node[sc] (s2) at (5.6,-1.6) {scene: \texttt{WindowGroup(for:)}};
  \draw[flow] (app.south -| s1) -- (s1);
  \draw[flow] (6.75,0 |- app.south) -- (6.75,0 |- s2.north);
  % windows
  \node[win] (w1) at (0.3,-2.75) {window 1\\\texttt{@State}, path,\\\texttt{@SceneStorage}};
  \node[win] (w2) at (2.6,-2.75) {window 2\\own copies};
  \node[win] (w3) at (5.6,-2.75) {note 42\\\texttt{\$id} = 42};
  \draw[flow] (s1) -- (w1); \draw[flow] (s1) -- (w2); \draw[flow] (s2) -- (w3);
  \node[nd, draw=sheetRed, fill=sheetRed!6, font=\tiny] (url) at (4.05,-2.2) {URL};
  \draw[hot] (url) -- (w3.west);
  % phases
  \node[ph, fill=sheetGreen!25] at (0.3,-3.55) {.active};
  \node[ph, fill=sheetOrange!25] at (2.6,-3.55) {.inactive};
  \node[ph, fill=black!12] at (5.6,-3.55) {.background};
  \draw[hot, rounded corners=3pt] (0.3,-3.75) -- (0.3,-4.05) -- (5.6,-4.05) -- (5.6,-3.75);
  \draw[hot] (2.6,-4.05) -- (2.6,-3.75);
  \node[lbl, anchor=north, text=sheetOrange, align=center] at (2.95,-4.1)
       {\texttt{App}'s \texttt{scenePhase} = \textbf{.active} (any window active) — a view's = its own window};
  \node[lbl, anchor=west, align=left, text=black!70, text width=78mm] at (-0.7,-4.8)
       {\texttt{store} in the App = shared; in a root view's \texttt{@State} = per window.
        {\color{sheetRed}URL}: \texttt{handlesExternalEvents} picks the window.};
\end{tikzpicture}

\section{Example}
\begin{lstlisting}[language=SwiftSheet]
@main struct NotesApp: App {
  @UIApplicationDelegateAdaptor(AppDelegate.self) var delegate
  @State private var store = NoteStore()     // shared
  @Environment(\.scenePhase) private var phase  // aggregate
  var body: some Scene {
    WindowGroup { RootView() }.environment(store)
      .onChange(of: phase) { _, p in
        if p == .background { store.save(); scheduleRefresh() } }
      .backgroundTask(.appRefresh("com.example.sync")) {
        await store.sync() }                    // return = done
    WindowGroup("Note", id: "note", for: Note.ID.self) { $id in
      NoteWindow(id: id) }.environment(store)
  } }
struct RootView: View {
  @SceneStorage("sel") private var sel: String?  // per window
  var body: some View { NoteList(sel: $sel)
    .onOpenURL { sel = Route($0)?.noteID } } }
\end{lstlisting}

\section{Traps}
\begin{itemize}
  \trap{Saving on \texttt{.inactive} only, or in one window's \texttt{onChange} — use the App-level
        \texttt{.background}; per-window phases fire once per window.}
  \trap{\texttt{@AppStorage} for a selection: two iPad windows overwrite each other — \texttt{@SceneStorage}.}
  \trap{\texttt{.backgroundTask} with no submitted request / Info.plist id never runs; timing is the system's.}
  \trap{\texttt{Window}, \texttt{Settings}, \texttt{MenuBarExtra} are not iOS — \texttt{\#if os(macOS)}.}
  \trap{\textbf{TN3187}: a UIKit app without the UIScene life cycle, built with the SDK after iOS 26,
        \textbf{won't launch}. SwiftUI \texttt{App} is scene-based already.}
\end{itemize}

\section{UIKit callback $\to$ SwiftUI}
{\footnotesize
\begin{tabular}{@{}L{37mm}L{40mm}@{}}
\toprule
\texttt{didFinishLaunching} & \texttt{App.init} / adaptor \\
\texttt{sceneDidBecomeActive} \ldots & \texttt{onChange(of: scenePhase)} \\
\texttt{scene(\_:openURLContexts:)} & \texttt{.onOpenURL} \\
\texttt{scene(\_:continue:)} & \texttt{.onContinueUserActivity} \\
\texttt{stateRestorationActivity} & \texttt{@SceneStorage}, \texttt{path.codable} \\
\texttt{BGTaskScheduler.register} & \texttt{.backgroundTask(.appRefresh)} \\
push token, \texttt{UNUserNotificationCenter} delegate & the adaptor's delegate methods \\
\bottomrule
\end{tabular}\par}

\end{multicols}

\noindent{\footnotesize\color{sheetGrey}\textit{Related:} state-restoration-multiwindow (UIKit scene model) ·
app-and-vc-lifecycle · background-execution · swiftui-navigation · deep-links-universal-links · push-notifications}

\end{document}
